% Options for packages loaded elsewhere
\PassOptionsToPackage{unicode}{hyperref}
\PassOptionsToPackage{hyphens}{url}
\documentclass[
]{article}
\usepackage{xcolor}
\usepackage[margin=1in]{geometry}
\usepackage{amsmath,amssymb}
\setcounter{secnumdepth}{5}
\usepackage{iftex}
\ifPDFTeX
  \usepackage[T1]{fontenc}
  \usepackage[utf8]{inputenc}
  \usepackage{textcomp} % provide euro and other symbols
\else % if luatex or xetex
  \usepackage{unicode-math} % this also loads fontspec
  \defaultfontfeatures{Scale=MatchLowercase}
  \defaultfontfeatures[\rmfamily]{Ligatures=TeX,Scale=1}
\fi
\usepackage{lmodern}
\ifPDFTeX\else
  % xetex/luatex font selection
\fi
% Use upquote if available, for straight quotes in verbatim environments
\IfFileExists{upquote.sty}{\usepackage{upquote}}{}
\IfFileExists{microtype.sty}{% use microtype if available
  \usepackage[]{microtype}
  \UseMicrotypeSet[protrusion]{basicmath} % disable protrusion for tt fonts
}{}
\makeatletter
\@ifundefined{KOMAClassName}{% if non-KOMA class
  \IfFileExists{parskip.sty}{%
    \usepackage{parskip}
  }{% else
    \setlength{\parindent}{0pt}
    \setlength{\parskip}{6pt plus 2pt minus 1pt}}
}{% if KOMA class
  \KOMAoptions{parskip=half}}
\makeatother
\usepackage{longtable,booktabs,array}
\newcounter{none} % for unnumbered tables
\usepackage{calc} % for calculating minipage widths
% Correct order of tables after \paragraph or \subparagraph
\usepackage{etoolbox}
\makeatletter
\patchcmd\longtable{\par}{\if@noskipsec\mbox{}\fi\par}{}{}
\makeatother
% Allow footnotes in longtable head/foot
\IfFileExists{footnotehyper.sty}{\usepackage{footnotehyper}}{\usepackage{footnote}}
\makesavenoteenv{longtable}
\usepackage{graphicx}
\makeatletter
\newsavebox\pandoc@box
\newcommand*\pandocbounded[1]{% scales image to fit in text height/width
  \sbox\pandoc@box{#1}%
  \Gscale@div\@tempa{\textheight}{\dimexpr\ht\pandoc@box+\dp\pandoc@box\relax}%
  \Gscale@div\@tempb{\linewidth}{\wd\pandoc@box}%
  \ifdim\@tempb\p@<\@tempa\p@\let\@tempa\@tempb\fi% select the smaller of both
  \ifdim\@tempa\p@<\p@\scalebox{\@tempa}{\usebox\pandoc@box}%
  \else\usebox{\pandoc@box}%
  \fi%
}
% Set default figure placement to htbp
\def\fps@figure{htbp}
\makeatother
\setlength{\emergencystretch}{3em} % prevent overfull lines
\providecommand{\tightlist}{%
  \setlength{\itemsep}{0pt}\setlength{\parskip}{0pt}}
\usepackage{amsmath}
\usepackage{booktabs}
\usepackage{array}
\usepackage{bookmark}
\IfFileExists{xurl.sty}{\usepackage{xurl}}{} % add URL line breaks if available
\urlstyle{same}
\hypersetup{
  pdftitle={ECON 3003 -- Homework \#1 Solutions},
  pdfauthor={Sarah},
  hidelinks,
  pdfcreator={LaTeX via pandoc}}

\title{ECON 3003 -- Homework \#1 Solutions}
\author{Sarah}
\date{February 27, 2026}

\begin{document}
\maketitle

\newpage

\section{Question 1 {[}7 pts{]} -- Extensive Form
Game}\label{question-1-7-pts-extensive-form-game}

\textbf{Strategic situation:} Javonie guesses which door Lycenia is
behind; Malik (host) chooses a color; Javonie picks a door.

\textbf{Players:} Lycenia (L), Malik (M), Javonie (J)

\textbf{Order of play:}

\begin{enumerate}
\def\labelenumi{\arabic{enumi}.}
\tightlist
\item
  Lycenia chooses A or B (Javonie not present)
\item
  Malik observes Lycenia's choice and chooses Red or Green
\item
  Javonie observes Malik's color (but not Lycenia's choice) and chooses
  A or B
\end{enumerate}

\textbf{Payoffs:} (Javonie, Lycenia, Malik)

\begin{itemize}
\tightlist
\item
  Javonie: \$100 if correct door, \$0 if wrong
\item
  Lycenia: \$100 if Javonie wrong, \$0 if correct
\item
  Malik: 10 if Javonie picks A, 0 if Javonie picks B
\end{itemize}

\textbf{Extensive form (decision tree):}

\begin{verbatim}
                    Lycenia
                   /      \
                 A          B
                /            \
              Malik          Malik
             /    \          /    \
          Red    Green    Red   Green
           |       |       |      |
        [cloud] [cloud] [cloud][cloud]  <- Javonie's information sets
           |       |       |      |
          Jav     Jav     Jav    Jav
         /  \    /  \    /  \   /  \
        A   B   A   B   A   B  A   B
\end{verbatim}

~~~~~~\textbar{} \textbar{} Jav Jav Jav Jav / ~ / ~ / ~ / ~ A B A B A B
A B

\textbackslash end\{verbatim\}

\textbf{Terminal payoffs (J, L, M):}

\begin{itemize}
\tightlist
\item
  Lycenia A, Malik Red: J picks A \(\rightarrow\) (100, 0, 10); J picks
  B \(\rightarrow\) (0, 100, 0)
\item
  Lycenia A, Malik Green: J picks A \(\rightarrow\) (100, 0, 10); J
  picks B \(\rightarrow\) (0, 100, 0)
\item
  Lycenia B, Malik Red: J picks A \(\rightarrow\) (0, 100, 10); J picks
  B \(\rightarrow\) (100, 0, 0)
\item
  Lycenia B, Malik Green: J picks A \(\rightarrow\) (0, 100, 10); J
  picks B \(\rightarrow\) (100, 0, 0)
\end{itemize}

\textbf{Information sets:}

\begin{itemize}
\tightlist
\item
  Lycenia: one decision node (no information set)
\item
  Malik: two nodes (knows Lycenia's choice)
\item
  Javonie: two information sets --- one when Malik chose Red, one when
  Malik chose Green (she does not know Lycenia's choice)
\end{itemize}

\newpage

\section{Question 2 {[}8 pts{]} -- Normal Form from Extensive
Form}\label{question-2-8-pts-normal-form-from-extensive-form}

\emph{Note: The extensive-form diagram is in the PDF. The method is:}

\begin{enumerate}
\def\labelenumi{\arabic{enumi}.}
\tightlist
\item
  Identify each player's strategies (complete plans of action for every
  node/information set)
\item
  Build the normal-form matrix with strategies as rows/columns
\item
  Enter payoffs for each strategy profile
\end{enumerate}

\textbf{General procedure:}

\begin{itemize}
\tightlist
\item
  For each information set, list the actions
\item
  A strategy for a player is a function from each of their information
  sets to an action
\item
  Enumerate all such strategies and fill the payoff matrix
\end{itemize}

\emph{(Apply this to the specific tree in the assignment.)}

\newpage

\section{Question 3 {[}6 pts total{]}}\label{question-3-6-pts-total}

\textbf{Game:}

{\def\LTcaptype{none} % do not increment counter
\begin{longtable}[]{@{}lcc@{}}
\toprule\noalign{}
& L & R \\
\midrule\noalign{}
\endhead
\bottomrule\noalign{}
\endlastfoot
U & 3, 3 & 2, 0 \\
D & 4, 1 & 8, -1 \\
Payoff & (P1, P2) & \\
\end{longtable}
}

\subsection{(i) Strictly and weakly dominated strategies
{[}2{]}}\label{i-strictly-and-weakly-dominated-strategies-2}

\textbf{Player 2:} - If P1 plays U: L gives 3, R gives 0 \(\rightarrow\)
L \(>\) R - If P1 plays D: L gives 1, R gives \(-1\) \(\rightarrow\) L
\(>\) R

So \textbf{R is strictly dominated by L} for Player 2.

\textbf{Player 1:} - If P2 plays L: U gives 3, D gives 4 \(\rightarrow\)
D \(>\) U - If P2 plays R: U gives 2, D gives 8 \(\rightarrow\) D \(>\)
U

So \textbf{U is strictly dominated by D} for Player 1.

\textbf{Weakly dominated:} None.

\subsection{(ii) Pure strategy Nash equilibrium
{[}2{]}}\label{ii-pure-strategy-nash-equilibrium-2}

After removing dominated strategies, only \textbf{(D, L)} remains.

\textbf{Check:} - Given P2 plays L, P1's best reply is D (4 \(>\) 3) -
Given P1 plays D, P2's best reply is L (1 \(>\) \(-1\))

\textbf{Nash equilibrium: (D, L) with payoff (4, 1).}

\subsection{(iii) Dominant strategy equilibrium?
{[}2{]}}\label{iii-dominant-strategy-equilibrium-2}

Yes. \textbf{D} is a dominant strategy for Player 1, and \textbf{L} is a
dominant strategy for Player 2. So \textbf{(D, L)} is a \textbf{dominant
strategy equilibrium} as well as Nash.

\newpage

\section{Question 4 {[}8 pts total{]}}\label{question-4-8-pts-total}

\textbf{Game:}

{\def\LTcaptype{none} % do not increment counter
\begin{longtable}[]{@{}lccc@{}}
\toprule\noalign{}
& L & C & R \\
\midrule\noalign{}
\endhead
\bottomrule\noalign{}
\endlastfoot
U & 5, 9 & 0, 1 & 4, 3 \\
M & 3, 2 & 0, 9 & 1, 1 \\
D & 2, 8 & 0, 1 & 8, 4 \\
Payoff & (P1, P2) & & \\
\end{longtable}
}

\subsection{(i) Weakly and strictly dominated strategies
{[}2{]}}\label{i-weakly-and-strictly-dominated-strategies-2}

\textbf{Player 2:} R is strictly dominated by L (L \(\geq\) R in every
row, strict when P1 plays U or D).

\textbf{Player 1:} No strategy strictly dominated in the full game.
After removing R, M and D are strictly dominated by U when P2 plays L.

\subsection{(ii) Pure strategy Nash equilibria
{[}4{]}}\label{ii-pure-strategy-nash-equilibria-4}

\textbf{Best responses:} - BR\(_1\)(L) = U; BR\(_1\)(C) = U, M, D;
BR\(_1\)(R) = D - BR\(_2\)(U) = L; BR\(_2\)(M) = C; BR\(_2\)(D) = L

\textbf{Mutual best replies:} - (U, L): BR\(_1\)(L) = U, BR\(_2\)(U) = L
\(\checkmark\) - (M, C): BR\(_1\)(C) includes M, BR\(_2\)(M) = C
\(\checkmark\)

\textbf{Nash equilibria: (U, L) with payoff (5, 9) and (M, C) with
payoff (0, 9).}

\subsection{(iii) Iterated dominance
{[}2{]}}\label{iii-iterated-dominance-2}

\textbf{Iterated elimination:} 1. R strictly dominated by L for P2
\(\rightarrow\) remove R 2. M and D strictly dominated by U for P1
\(\rightarrow\) remove M and D 3. P2 prefers L (9 \(>\) 1)

\textbf{Iterated dominance equilibrium: (U, L) with payoff (5, 9).}

So \textbf{(U, L)} can be reached by iterated dominance; \textbf{(M, C)}
cannot.

\newpage

\section{Question 5 {[}16 pts total{]}}\label{question-5-16-pts-total}

\textbf{Game:}

{\def\LTcaptype{none} % do not increment counter
\begin{longtable}[]{@{}lcc@{}}
\toprule\noalign{}
& L & R \\
\midrule\noalign{}
\endhead
\bottomrule\noalign{}
\endlastfoot
U & 3, 2 & 2, 1 \\
D & 4, 3 & 1, 4 \\
Payoff & (P1, P2) & \\
\end{longtable}
}

\subsection{(i) Pure strategy Nash equilibrium?
{[}4{]}}\label{i-pure-strategy-nash-equilibrium-4}

\textbf{Best responses:} - BR\(_1\)(L) = D; BR\(_1\)(R) = U -
BR\(_2\)(U) = L; BR\(_2\)(D) = R

No profile is a mutual best reply. \textbf{No Nash equilibrium in pure
strategies.}

\subsection{(ii) Mixed strategy Nash equilibrium
{[}10{]}}\label{ii-mixed-strategy-nash-equilibrium-10}

Let P1 play U with probability \(p\), D with \(1-p\). Let P2 play L with
probability \(q\), R with \(1-q\).

\textbf{P2 indifference:} \[
2p + 3(1-p) = 1p + 4(1-p) \Rightarrow -p + 3 = -3p + 4 \Rightarrow 2p = 1 \Rightarrow p^* = \frac{1}{2}
\]

\textbf{P1 indifference:} \[
3q + 2(1-q) = 4q + 1(1-q) \Rightarrow q + 2 = 3q + 1 \Rightarrow 2q = 1 \Rightarrow q^* = \frac{1}{2}
\]

\textbf{Mixed strategy Nash equilibrium:} - P1:
\((\tfrac{1}{2} U, \tfrac{1}{2} D)\) - P2:
\((\tfrac{1}{2} L, \tfrac{1}{2} R)\)

\textbf{Expected payoffs:} (2.5, 2.5)

\subsection{(iii) Pareto optimality
{[}2{]}}\label{iii-pareto-optimality-2}

Pure outcomes: (3,2), (2,1), (4,3), (1,4). The profile (4,3) Pareto
dominates (2.5, 2.5). So the mixed strategy Nash equilibrium is
\textbf{not Pareto optimal}.

\end{document}
